\section{Related Work}
\label{sec:rw}

Static analysis tools for Infrastructure-as-Code (IaC) such as Checkov and tfsec catch a well-known set of syntactically enumerable misconfigurations (permissive file modes, unencrypted protocols, world-open network rules) reliably, but struggle with misconfigurations that depend on context or intent rather than a fixed pattern.

\subsection{Baseline: Semantics-Aware Detection with Code and Language Models}
The baseline for this project is War et al. \cite{War25}, ``Detection of Security Smells in IaC Scripts through Semantics-Aware Code and Language Processing'' (University of Luxembourg, arXiv:2509.18790, September 2025). \emph{This is cited as a preprint; its peer-review status could not be confirmed at the time of writing, unlike the baselines used elsewhere in this thesis series.} The paper fine-tunes CodeBERT (for Ansible scripts, which are rich in natural-language comments and task descriptions) and LongFormer (for long Puppet scripts) to classify IaC snippets as misconfigured or safe, substantially outperforming prior TF-IDF/Bag-of-Words + Random Forest baselines (e.g. Ansible F1 0.90 vs. 0.77) and performing competitively with much larger general-purpose LLMs.

\subsection{The Gap This Project Targets}
The paper's own ablation study (their RQ2) is the source of the gap this project targets. Removing natural-language context -- stripping comments from Ansible snippets, or truncating surrounding code for Puppet snippets -- causes precision to collapse sharply while recall holds steady or rises: CodeBERT on Ansible drops from precision 0.92 to 0.46 (F1 0.90 to 0.58); LongFormer on Puppet drops from 0.87 to 0.55 (F1 0.79 to 0.70). Their own words: ``without sufficient semantic context... models over-flag and lose precision, undermining practical usability.'' This is a real deployment concern: much real-world IaC is sparsely commented, unlike the curated, comment-rich datasets used to train and validate the paper's models.

\subsection{Hybrid Rule-Based and ML Detection}
Combining rule-based static analysis with ML classifiers is a natural response to this specific failure mode -- a fixed rule cannot lose precision due to missing comments, because it never used comments to begin with. This project builds and evaluates exactly that combination on a synthetic benchmark designed so that some misconfigurations are enumerable (rule-detectable) and some are only resolvable via natural-language context (structurally rule-proof), to test how much of the gap a hybrid approach can close, and to make explicit what it structurally cannot.
